\pagestyle{plain}
\songti
\begin{center}
    \Large
    \textbf{离散数学复习重点}
\end{center}


\noindent\textbf{一、真值表法}

析取 \(\lor\) 只要不是 \(00\) 就是 \(1\)，
合取 \(\land\) 只要不是 \(11\) 就是 \(0\)，
\(\to\) 只要不是 \(10\) 就是 \(1\)，
\(\leftrightarrow\) 只要不是 \(11\) 和 \(00\) 就是 \(0\)，
\(\neg\) 取反。

值为 $1$ 的真值判定范畴称为\textbf{小真}，用小写字母 $m$ 加下标表示，二进制用析取连接；值为 $0$ 的真值判定范畴称为\textbf{大真}，用大写字母 $M$ 加下标表示，二进制用合取连接。

\noindent\textbf{二、等值演算}

\begin{enumerate}
    \item 双重否定律
    \begin{align*}
        \neg \neg P &\Leftrightarrow P
    \end{align*}

    \item 等幂律
    \begin{align*}
        P \vee P &\Leftrightarrow P, &
        P \wedge P &\Leftrightarrow P
    \end{align*}

    \item 交换律
    \begin{align*}
        P \vee Q &\Leftrightarrow Q \vee P, &
        P \wedge Q &\Leftrightarrow Q \wedge P
    \end{align*}

    \item 结合律
    \begin{align*}
        P \vee (Q \vee R) &\Leftrightarrow (P \vee Q) \vee R, \\
        P \wedge (Q \wedge R) &\Leftrightarrow (P \wedge Q) \wedge R
    \end{align*}

    \item 分配律
    \begin{align*}
        P \vee (Q \wedge R) &\Leftrightarrow (P \vee Q) \wedge (P \vee R), \\
        P \wedge (Q \vee R) &\Leftrightarrow (P \wedge Q) \vee (P \wedge R)
    \end{align*}

    \item 德摩根定律
    \begin{align*}
        \neg (P \vee Q) &\Leftrightarrow \neg P \wedge \neg Q, \\
        \neg (P \wedge Q) &\Leftrightarrow \neg P \vee \neg Q
    \end{align*}

    \item 吸收律
    \begin{align*}
        P \vee (P \wedge Q) &\Leftrightarrow P, &
        P \wedge (P \vee Q) &\Leftrightarrow P
    \end{align*}

    \item 蕴含等值式
    \begin{align*}
        P \to Q &\Leftrightarrow \neg P \vee Q
    \end{align*}

    \item 等价等值式
    \begin{align*}
        P \leftrightarrow Q &\Leftrightarrow (P \to Q) \wedge (Q \to P)
    \end{align*}

    \item 零律
    \begin{align*}
        P \vee T &\Leftrightarrow T, &
        P \wedge F &\Leftrightarrow F
    \end{align*}

    \item 同一律
    \begin{align*}
        P \vee F &\Leftrightarrow P, &
        P \wedge T &\Leftrightarrow P
    \end{align*}

    \item 排中律
    \begin{align*}
        P \vee \neg P &\Leftrightarrow T
    \end{align*}

    \item 矛盾律
    \begin{align*}
        P \wedge \neg P &\Leftrightarrow F
    \end{align*}

    \item 等价否定等值式
    \begin{align*}
        P \leftrightarrow Q &\Leftrightarrow (\neg P \leftrightarrow \neg Q)
    \end{align*}

    \item 归谬律
    \begin{align*}
        (P \to Q) \wedge (P \to \neg Q) &\Leftrightarrow \neg P
    \end{align*}

    \item 假言移位
    \begin{align*}
        P \to Q &\Leftrightarrow \neg Q \to \neg P
    \end{align*}
\end{enumerate}

\noindent\textbf{三、推演定律}


\begin{table}[H]
    \centering
    \begin{tabular}{cc}
        \toprule
        \textbf{公式} & \textbf{律名} \\
        \midrule
        $A \Rightarrow (A \lor B)$ & 附加律 \\
        $(A \land B) \Rightarrow A$ & 化简律 \\
        $(A \to B) \land A \Rightarrow B$ & 假言推理 \\
        $(A \to B) \land \neg B \Rightarrow \neg A$ & 拒取式 \\
        $(A \lor B) \land \neg B \Rightarrow A$ & 析取三段论 \\
        \bottomrule
    \end{tabular}
\end{table}

\noindent\textbf{四、前束范式}

\begin{enumerate}
    \item 消去除 $\neg$、$\wedge$、$\vee$ 之外的联结词；
    \item 将否定符 $\neg$ 移到量词符后；
    \item 换名使命题变元不同名；
    \item 扩大辖域使所有量词处在最前面。
\end{enumerate}

\noindent\textbf{五、闭包运算}

设二元关系为 $R\subseteq A\times A$。

\begin{enumerate}
    \item 自反：对任意 $a\in A$，都有 $(a,a)\in R$。例：$\le$ 在整数集上自反，因为 $a\le a$。


    \item 反自反：对任意 $a\in A$，都有 $(a,a)\notin R$。例：$<$ 在整数集上反自反，因为不存在 $a<a$。


    \item 对称：若 $(a,b)\in R$，则 $(b,a)\in R$。例：“同学关系”通常是对称的，若 $a$ 是 $b$ 的同学，则 $b$ 也是 $a$ 的同学。


    \item 反对称：若 $(a,b)\in R$ 且 $(b,a)\in R$，则 $a=b$。例：$\le$ 是反对称的，因为若 $a\le b$ 且 $b\le a$，则 $a=b$。


    \item 传递：若 $(a,b)\in R$ 且 $(b,c)\in R$，则 $(a,c)\in R$。例：$<$ 是传递的，因为若 $a<b$ 且 $b<c$，则 $a<c$。


    \item 非传递：存在 $a,b,c\in A$，使得 $(a,b)\in R$ 且 $(b,c)\in R$，但 $(a,c)\notin R$。例：“朋友关系”一般非传递，若 $a$ 是 $b$ 的朋友，$b$ 是 $c$ 的朋友，不一定 $a$ 是 $c$ 的朋友。

\end{enumerate}



\noindent\textbf{例题:} 设 $A=\{1,2,3,4\}$，$R=\{(1,2),(2,4),(3,3),(1,3)\}$。

\begin{enumerate}
    \item 自反闭包：因为自反关系要求对任意 $a\in A$，都有 $(a,a)\in R$。
    
    原关系中已有 $(3,3)$，还需要补充 $(1,1)$、$(2,2)$、$(4,4)$.
    
    所以 $r(R)=\{(1,2),(2,4),(3,3),(1,3),(1,1),(2,2),(4,4)\}$。
    
    \item 对称闭包：因为对称关系要求若 $(a,b)\in R$，则 $(b,a)\in R$。
    
    由 $(1,2)$ 需补充 $(2,1)$，由 $(2,4)$ 需补充 $(4,2)$，由 $(1,3)$ 需补充 $(3,1)$，而 $(3,3)$ 不需要补充，所以 $s(R)=\{(1,2),(2,4),(3,3),(1,3),(2,1),(4,2),(3,1)\}$。
    
    \item 传递闭包：因为传递关系要求若 $(a,b)\in R$ 且 $(b,c)\in R$，则 $(a,c)\in R$。
    
    由 $(1,2)\in R$ 且 $(2,4)\in R$，需补充 $(1,4)$。其他有序对不能继续推出新的有序对，所以 $t(R)=\{(1,2),(2,4),(3,3),(1,3),(1,4)\}$。
\end{enumerate}


\noindent\textbf{六、关系矩阵}

设集合 \(A=\{1,2,3,4\}\)，关系 \(R=\{(1,3),(1,4),(2,3),(2,4),(3,4)\}\)。

\begin{enumerate}
    \item \(R\) 的关系图为：顶点为 \(1,2,3,4\)，有向边为 \(1\to3\)，\(1\to4\)，\(2\to3\)，\(2\to4\)，\(3\to4\)。

    \item 按 \(A\) 中元素顺序 \(1,2,3,4\) 排列，\(R\) 的关系矩阵为 \(M_R=\begin{pmatrix}0&0&1&1\\0&0&1&1\\0&0&0&1\\0&0&0&0\end{pmatrix}\)。

    \item \(R\) 不具有自反性，因为 \((1,1)\notin R\)，\((2,2)\notin R\)，\((3,3)\notin R\)，\((4,4)\notin R\)。\(R\) 具有反自反性，因为对任意 \(a\in A\)，都有 \((a,a)\notin R\)。\(R\) 不具有对称性，因为 \((1,3)\in R\)，但 \((3,1)\notin R\)。\(R\) 具有反对称性，因为不存在 \(a,b\in A\) 且 \(a\ne b\)，使得 \((a,b)\in R\) 且 \((b,a)\in R\)。
\end{enumerate}

\noindent\textbf{七、复合运算}

设$
R = \{\langle 1,4\rangle,\langle 3,5\rangle,\langle 2,5\rangle\},S = \{\langle 4,2\rangle,\langle 5,1\rangle,\langle 2,3\rangle\}$是集合$A=\{1,2,3,4,5\}$上的两个关系，求$R\circ S,S\circ R$。

\textbf{解答}

按关系复合“先左后右”的记法：
$$R\cdot S=\{\langle x,z\rangle \mid \exists y,\ \langle x,y\rangle\in R,\ \langle y,z\rangle\in S\}$$

逐个算：

$\langle1,4\rangle\in R,\ \langle4,2\rangle\in S \Rightarrow \langle1,2\rangle$
$\langle3,5\rangle\in R,\ \langle5,1\rangle\in S \Rightarrow \langle3,1\rangle$
$\langle2,5\rangle\in R,\ \langle5,1\rangle\in S \Rightarrow \langle2,1\rangle$

所以：
$$R\cdot S=\{\langle1,2\rangle,\langle3,1\rangle,\langle2,1\rangle\}$$

同理：

$\langle4,2\rangle\in S,\ \langle2,5\rangle\in R \Rightarrow \langle4,5\rangle$
$\langle5,1\rangle\in S,\ \langle1,4\rangle\in R \Rightarrow \langle5,4\rangle$
$\langle2,3\rangle\in S,\ \langle3,5\rangle\in R \Rightarrow \langle2,5\rangle$

因此：
$$S\cdot R=\{\langle4,5\rangle,\langle5,4\rangle,\langle2,5\rangle\}$$

答案：
$$\boxed{R\cdot S=\{\langle1,2\rangle,\langle3,1\rangle,\langle2,1\rangle\}}$$
$$\boxed{S\cdot R=\{\langle4,5\rangle,\langle5,4\rangle,\langle2,5\rangle\}}$$